\clearpage\fancyhead[L]{\small VESTA.}\hypertarget{native-382-lower}{}\label{native-382-lower}
\begin{center}\textbf{VESTA.}\par\vspace{4pt}\hrule\vspace{4pt}\fontsize{10}{12}\selectfont Göttingische gelehrte Anzeigen. Stück 129. Seite 1281\;..\;1284. 1814. August 13.\par\vspace{4pt}\hrule\vspace{4pt}\end{center}

\par\indent Wir haben noch eine kurze Anzeige von der letzten Opposition der \emph{Vesta} im Februar d. J.
und einigen dadurch veranlassten Rechnungen nachzuholen. Auf der hiesigen Sternwarte waren von
dem Prof. \textsc{Gauss} folgende Beobachtungen am Mauerquadranten gemacht worden:

\begin{center}\fontsize{9}{11}\selectfont\setlength{\tabcolsep}{3pt}\renewcommand{\arraystretch}{1.15}\begin{tabular}{l|c|rrr|rrr}
1814 & \shortstack{Mittlere Zeit\\in Göttingen} & \multicolumn{3}{c|}{\shortstack{Scheinbare gerade\\Aufsteigung}} & \multicolumn{3}{c}{\shortstack{Scheinbare\\Abweichung}}\\\hline
Febr. 12 & $12^{\mathrm u}\;30^{\mathrm m}\;15^{\mathrm s}$ & $149^\circ$ & $56'$ & $48''3$ & $20^\circ$ & $47'$ & $3''3$ N.\\
14 & $12\;20\;26$ & 149 & 27 & 22.8 & 21 & 3 & 3.3\\
15 & $12\;15\;31$ & 149 & 12 & 35.5 & 21 & 10 & 55.7\\
20 & $11\;50\;55$ & 147 & 58 & 14.3 & 21 & 47 & 35.5\\
\end{tabular}\end{center}

{\rightskip=0pt plus 1em \par\indent Bei der sehr geringen Beleuchtung, die der Planet in dem lichtschwachen Fernrohr vertrug, kann
diesen Beobachtungen nur eine mittelmässige Genauigkeit beigelegt werden.\par}

\SourceLeaf{383}\Needspace{20\baselineskip}\par\indent Von auswärtigen Sternwarten erhielt Prof. \textsc{Gauss} folgende Beobachtungen mitgetheilt:

\Needspace{14\baselineskip}\begin{center}\emph{Beobachtungen der Vesta auf der Seeberger Sternwarte von Hrn.} \textsc{Nicolai}.\end{center}
\begin{center}\fontsize{9}{11}\selectfont\setlength{\tabcolsep}{3pt}\renewcommand{\arraystretch}{1.15}\begin{tabular}{l|c|rrr|c}
1814 & \shortstack{Mittlere Seeberger\\Zeit} & \multicolumn{3}{c|}{\shortstack{Scheinbare gerade\\Aufsteigung}} & \shortstack{Scheinbare\\Abweichung}\\\hline
Febr. 12 & $12^{\mathrm u}\;30^{\mathrm m}\;16^{\mathrm s}0$ & $149^\circ$ & $57'$ & $0''0$ & $\cdot\quad\cdot\quad\cdot$\\
15 & $12\;15\;31.6$ & 149 & 12 & 42.0 & $\cdot\quad\cdot\quad\cdot$\\
20 & $11\;50\;55.1$ & 147 & 58 & 15.9 & $21^\circ\;47'\;28''$ N.\\
25 & $11\;26\;25.1$ & 146 & 45 & 26.8 & $\cdot\quad\cdot\quad\cdot$\\
26 & $11\;21\;32.7$ & 146 & 31 & 17.7 & $\cdot\quad\cdot\quad\cdot$\\
28 & $11\;11\;50.7$ & 146 & 3 & 39.9 & $\cdot\quad\cdot\quad\cdot$\\
\end{tabular}\end{center}
\Needspace{17\baselineskip}\begin{center}\emph{Beobachtungen der Vesta auf der Königsberger Sternwarte von Hrn. Prof.} \textsc{Bessel}.\end{center}
\begin{center}\fontsize{9}{11}\selectfont\setlength{\tabcolsep}{3pt}\renewcommand{\arraystretch}{1.15}\begin{tabular}{l|rrr|rrr|rrr}
1814 & \multicolumn{3}{c|}{\shortstack{Mittlere Zeit\\in Königsberg}} & \multicolumn{3}{c|}{\shortstack{Scheinbare gerade\\Aufsteigung}} & \multicolumn{3}{c}{\shortstack{Scheinbare\\Abweichung}} \\
\hline
Jan.\ 31  & $13^{\mathrm u}$ & $28^{\mathrm m}$ & $25^{\mathrm s}$ & $152^{\circ}$ & $40'$ & $16''8$ & $19^{\circ}$ & $6'$  & $38''4$ N. \\
Febr.\ 9  & 12 & 45 & 5  & 150 & 40 & 49.2 & 20 & 22 & 8.3 \\
13        & 12 & 25 & 30 & 149 & 42 & 44.3 & 20 & 54 & 38.2 \\
14        & 12 & 20 & 35 & 149 & 27 & 56.6 & 21 & 2  & 37.9 \\
19        & 11 & 55 & 58 & 148 & 13 & 33.3 & 21 & 40 & 19.8 \\
21        & 11 & 46 & 8  & 147 & 43 & 51.5 & 21 & 54 & 27.8 \\
22        & 11 & 41 & 13 & 147 & 29 & 9.1  & 22 & 1  & 9.7  \\
25        & 11 & 26 & 33 & 146 & 45 & 47.0 & 22 & 20 & 34.1 \\
26        & 11 & 21 & 40 & 146 & 31 & 36.8 & 22 & 26 & 34.1 \\
\end{tabular}\end{center}

\par\indent Die Polhöhe der neuen Königsberger Sternwarte hat Hr. Prof. \textsc{Bessel} aus 69 mit dem
2fussigen \textsc{Cary}'schen Kreise beobachteten Culminationen des Polarsterns zu
$54^{\circ}\,42'\,50''31$ bestimmt.

\Needspace{28\baselineskip}\par\indent Hr. Dr. \textsc{Gerling} in Cassel übernahm die Vergleichung dieser sämmtlichen Beobachtungen
mit den letzten Elementen, bei welchen, um eine bessere Uebereinstimmung zu erhalten, die
Epoche abgeändert und für den 13. Februar $12^{\mathrm u}$ M.\,Z.\ in Göttingen zu
$154^{\circ}\,54'\,8''39$ angenommen wurde. Diese Rechnung gab folgende Unterschiede

\begin{center}\fontsize{10}{12}\selectfont\setlength{\tabcolsep}{6pt}\renewcommand{\arraystretch}{1.15}\begin{tabular}{l|c|c|l}
& \shortstack{in gerader\\Aufsteigung} & \shortstack{in der\\Abweichung} & Beobachter \\
\hline
Jan.\ 31 & $-11''15$ & $+4''44$  & Bessel \\
Febr.\ 9 & $-6.05$   & $+5.08$   & Bessel \\
12       & $-6.75$   & $\cdot\;\cdot\;\cdot\;\cdot$ & Nicolai \\
12       & $-16.45$  & $+10.42$  & Gauss \\
13       & $-3.16$   & $-0.11$   & Bessel \\
14       & $-3.13$   & $+5.62$   & Bessel \\
14       & $-11.23$  & $+16.76$  & Gauss \\
15       & $-8.95$   & $+22.86$  & Gauss \\
15       & $-0.46$   & $\cdot\;\cdot\;\cdot\;\cdot$ & Nicolai \\
19       & $+4.77$   & $+4.42$   & Bessel \\
20       & $+2.53$   & $+6.57$   & Nicolai \\
20       & $+2.98$   & $+0.21$   & Gauss \\
21       & $+2.61$   & $+6.97$   & Bessel \\
22       & $+2.48$   & $+0.93$   & Bessel \\
25       & $+3.60$   & $+5.05$   & Bessel \\
25       & $+6.28$   & $\cdot\;\cdot\;\cdot\;\cdot$ & Nicolai \\
26       & $+1.76$   & $+0.66$   & Bessel \\
26       & $+5.58$   & $\cdot\;\cdot\;\cdot\;\cdot$ & Nicolai \\
28       & $+7.92$   & $\cdot\;\cdot\;\cdot\;\cdot$ & Nicolai \\
\end{tabular}\end{center}

\par\indent Hiernächst leitete Hr. Doctor \textsc{Gerling} aus diesen Resultaten die Opposition selbst
ab, und bestimmte dann neue Elemente, welche die vier Oppositionen von 1810, 1811, 1812,
1814 darstellen.

\Needspace{12\baselineskip}\SourceLeaf{384}\emph{Fünfte beobachtete Opposition der Vesta.}

\begin{center}1814 Februar 13. $9^{\mathrm u}\;5^{\mathrm m}\;40^{\mathrm s}$ Mittlere Zeit in Göttingen\end{center}
\begin{center}\fontsize{10}{12}\selectfont\setlength{\tabcolsep}{2pt}\renewcommand{\arraystretch}{1.15}\begin{tabular}{lcr}
Wahre Länge & $.\;.\;.\;.\;.\;.\;.\;.\;.\;.\;.\;.$ & $144^\circ\;34'\;54''47$\\
Geocentrische Breite & $.\;.\;.\;.\;.\;.\;.\;.\;.\;.$ & $8\;2\;8.55\;\mathrm{N.}$\\
\end{tabular}\end{center}

\Needspace{16\baselineskip}\emph{Neue Elemente der Vesta.}

\begin{center}\fontsize{9}{11}\selectfont\setlength{\tabcolsep}{2pt}\renewcommand{\arraystretch}{1.15}\begin{tabular}{lcr}
Epoche der mittlern Länge 1814 Febr. 13. & & $12^{\mathrm u}$\\
in Göttingen & & $154^\circ\;55'\;27''83$\\
Länge des Perihel & $.\;.\;.\;.\;.\;.\;.\;.\;.\;.\;.\;.\;.$ & $249\;38\;6.69$\\
Länge des aufsteigenden Knoten & $.\;.\;.\;.\;.\;.\;.\;.$ & $103\;11\;30.51$\\
\multicolumn{3}{c}{Beide für die Zeit der Epoche}\\
Neigung der Bahn & $.\;.\;.\;.\;.\;.\;.\;.\;.\;.\;.\;.\;.$ & $7\;8\;16.01$\\
Excentricitätswinkel & $.\;.\;.\;.\;.\;.\;.\;.\;.\;.\;.\;.$ & $5\;8\;30.75$\\
Logarithm. des mittlern Abstandes & $.\;.\;.\;.\;.\;.\;.$ & $0.3731261$\\
Mittlere tägliche tropische Bewegung & $.\;.\;.\;.\;.\;.$ & $977''95156$\\
\end{tabular}\end{center}

{\rightskip=0pt plus 1em \par\indent Derselbe geschickte Astronom ist jetzt mit der Berechnung einer Ephemeride für die nächste
Erscheinung der Vesta beschäftigt, welche demnächst an einem andern Orte bekannt gemacht
werden wird.\par}
